\section{The Canonical Ensemble}

\paragraph{Constraining the average of an observable.}
The microcanonical ensemble does not distinguish between accessible
microstates: only their number matters. Realistic systems, however, obey
constraints imposed by measurement. Suppose microstate $i$ has a definite
value $a_i$ of an observable $A$, and the measured average is $\bar A$. The
expectation value is
\begin{equation}
    \boxed{\langle A\rangle \equiv \sum_i p_i a_i,}
\end{equation}
and the constraint reads $\langle A\rangle=\bar A$.
\textcolor{red}{Here, $\langle A\rangle$ denotes the operation that extracts
the mean of the observable $A$, and $\bar A$ denotes the numerical value
assigned to that mean. Since only one constraint is present, we also rename
the multiplier $\lambda_a$ as $\beta$.}
The Lagrangian now contains the multiplier $\alpha$ for normalization and
the multiplier $\beta$ for the constraint on $\langle A\rangle$. Its
stationarity conditions are
\begin{align*}
    \frac{\partial\mathcal{L}}{\partial p_i}
    &= -\ln p_i-1-\alpha-\beta a_i=0, \\
    \frac{\partial\mathcal{L}}{\partial\alpha}
    &= -\Big(\sum_i p_i-1\Big)=0, \\
    \frac{\partial\mathcal{L}}{\partial\beta}
    &= -\Big(\sum_i a_i p_i-\bar A\Big)=0.
\end{align*}

\paragraph{The Boltzmann distribution and the partition function.}
The first condition gives $p_i = e^{-1-\alpha-\beta a_i}$. Substituting this
into normalization and into the constraint yields
\begin{align*}
    e^{-1-\alpha}\sum_i e^{-\beta a_i} = 1,
    \qquad
    e^{-1-\alpha}\sum_i a_i e^{-\beta a_i} = \bar A.
\end{align*}
The first equation determines $\alpha$ in terms of $\beta$ and leaves a
single free parameter. Defining the \textcolor{blue}{canonical partition
function}
\begin{equation}
    \boxed{Z(\beta) \equiv \sum_i e^{-\beta a_i},}
\end{equation}
the probability distribution becomes
\begin{equation}
    \boxed{p_i = \frac{e^{-\beta a_i}}{Z(\beta)}.}
\end{equation}
This is the \textcolor{blue}{Boltzmann distribution}. For $\beta>0$, states
with larger $a_i$ are less probable. The remaining constraint,
$\langle A\rangle(\beta)=\bar A$, determines $\beta$.

\paragraph{Uniqueness of the multiplier.}
To show that $\beta$ is unique, we first express the mean through $Z$:
\begin{align*}
    \langle A\rangle = \frac{1}{Z}\sum_i a_i e^{-\beta a_i}
    = -\frac{\partial\ln Z}{\partial\beta}.
\end{align*}
Hence
\begin{equation}
    \boxed{\langle A\rangle = -\frac{\partial\ln Z}{\partial\beta}.}
\end{equation}
Differentiating once more, with $\mathrm{Var}(A)\equiv\langle A^2\rangle-\langle A\rangle^2$,
\begin{align*}
    \frac{\partial\langle A\rangle}{\partial\beta}
    &= \frac{\partial}{\partial\beta}
    \left(\frac{\sum_i a_i e^{-\beta a_i}}{Z}\right) \\
    &= -\frac{\sum_i a_i^2 e^{-\beta a_i}}{Z}
    +\frac{1}{Z^2}\left(\sum_i a_i e^{-\beta a_i}\right)^2 \\
    &= -\langle A^2\rangle+\langle A\rangle^2
    = -\mathrm{Var}(A),
\end{align*}
so
\begin{equation}
    \boxed{\frac{\partial\langle A\rangle}{\partial\beta} = -\mathrm{Var}(A).}
\end{equation}
Since $\mathrm{Var}(A)\geq0$, $\langle A\rangle(\beta)$ is monotonically
non-increasing. It is strictly decreasing whenever $A$ is not identical
across all accessible microstates. Thus, for target values in the interior
of the achievable range, $\langle A\rangle$ determines $\beta$ uniquely.

\paragraph{Entropy of the canonical distribution and the meaning of $\beta$.}
We now compute the entropy of the resulting distribution, with $k_B$
omitted during the derivation:
\begin{align*}
    S &= -\sum_i p_i\ln p_i
    = -\sum_i p_i\left(-\beta a_i-\ln Z\right) \\
    &= \beta\sum_i a_i p_i+\ln Z\sum_i p_i
    = \beta\langle A\rangle+\ln Z.
\end{align*}
Restoring $k_B$,
\begin{equation}
    \boxed{S=k_B\left(\beta\langle A\rangle+\ln Z\right).}
\end{equation}
Although $S$ depends on $\bar A$ both explicitly and implicitly through
$\beta(\bar A)$, differentiating gives
\begin{align*}
    \frac{\partial S}{\partial\langle A\rangle}
    &= \beta
    +\langle A\rangle\frac{d\beta}{d\langle A\rangle}
    +\frac{\partial\ln Z}{\partial\beta}\frac{d\beta}{d\langle A\rangle} \\
    &= \beta
    +\langle A\rangle\frac{d\beta}{d\langle A\rangle}
    -\langle A\rangle\frac{d\beta}{d\langle A\rangle}
    = \beta.
\end{align*}
Restoring $k_B$ yields the meaning of the multiplier:
\begin{equation}
    \boxed{\beta=\frac{1}{k_B}\frac{\partial S}{\partial\langle A\rangle}.}
\end{equation}

\paragraph{Concavity of the entropy.}
A function $f$ is \textcolor{blue}{concave} on an interval $I$ if, for all
$x,y\in I$ and all $t\in[0,1]$,
\begin{equation}
    \boxed{f\big(tx+(1-t)y\big)\;\geq\;t\,f(x)+(1-t)\,f(y).}
\end{equation}
Geometrically, every chord lies on or below the graph. The function is
\textcolor{blue}{strictly concave} if the inequality is strict for $x\neq y$
and $t\in(0,1)$. For a twice-differentiable $f$, concavity is equivalent to
$f''\leq0$, and $f''<0$ guarantees strict concavity.

We apply this criterion to $S$ as a function of $\langle A\rangle$. Using
the inverse-function rule and $\partial\langle A\rangle/\partial\beta=-\mathrm{Var}(A)$,
\begin{align*}
    \frac{\partial^2 S}{\partial\langle A\rangle^2}
    = \frac{d\beta}{d\langle A\rangle}
    = \left(\frac{d\langle A\rangle}{d\beta}\right)^{-1}
    = -\frac{1}{\mathrm{Var}(A)}.
\end{align*}
Restoring $k_B$,
\begin{equation}
    \boxed{\frac{\partial^2 S}{\partial\langle A\rangle^2}
    = -\frac{k_B}{\mathrm{Var}(A)}<0.}
\end{equation}
The entropy is therefore strictly concave in $\langle A\rangle$ whenever $A$
is not identical across all accessible microstates. This is the monotonic
decrease of $\langle A\rangle(\beta)$ read from the side of $S$: each value
of $\langle A\rangle$ corresponds to exactly one $\beta$.

\paragraph{Specialization to the internal energy.}
We have not specified the observable $A$, so the derivation applies to any
observable whose expectation value is constrained. The canonical ensemble
conventionally fixes the volume $V$ and particle number $N$ and constrains
the internal energy. \textcolor{red}{From here on, we take $A=U$ and write
$U_i$ for $a_i$.} The results become
\begin{equation}
    \boxed{
    \begin{aligned}
        Z(\beta) &= \sum_i e^{-\beta U_i}, &
        p_i &= \frac{e^{-\beta U_i}}{Z(\beta)}, \\
        \langle U\rangle &= -\frac{\partial\ln Z}{\partial\beta}, &
        S &= k_B\left(\beta\langle U\rangle+\ln Z\right).
    \end{aligned}}
\end{equation}
We identify $\beta$ with the temperature in Section~\ref{sec:multipliers}.
The reason for not constraining only the particle number or the volume will
be discussed shortly.

\paragraph{Recovering the microcanonical ensemble.}
If every accessible microstate has the same energy $U$, the canonical
ensemble reduces to the microcanonical one. All weights $e^{-\beta U}$ are
identical, so the distribution is uniform and the mean energy is the exact
value $U$. Then
\begin{align*}
    Z &= \sum_{i=1}^{\Omega} e^{-\beta U} = \Omega\, e^{-\beta U}, \\
    S &= k_B\left(\beta U+\ln\Omega-\beta U\right) = k_B\ln\Omega.
\end{align*}

\subsection{First and Second Law of Thermodynamics - Making Sense of Multipliers}
\label{sec:multipliers}

\paragraph{The multipliers need thermodynamics for their meaning.}
The MaxEnt formalism does not by itself provide the physical interpretation
or numerical calibration of its multipliers. In particular, the
identification
\begin{equation}
    \boxed{\beta=\frac{1}{k_B}\frac{\partial S}{\partial \langle U\rangle}=\frac{1}{k_BT}}
\end{equation}
cannot be derived from MaxEnt alone. It is an empirical correspondence
established by classical thermodynamics. We therefore obtain the physical
meaning of the multipliers by connecting them to the thermodynamic laws.

\paragraph{From the first and second laws to the fundamental relation.}
The \textcolor{blue}{first law} expresses energy conservation:
\begin{equation}
    \boxed{dU = \delta Q-\delta W+\mu\, dN,}
\end{equation}
where
\begin{itemize}
    \item $\delta W=P\,dV$ is the work done \emph{by} the system.
    \item The symbols $\delta Q$ and $\delta W$ emphasize that heat and work
    are not state functions: their values depend on the path between two
    states.
    \item The \textcolor{blue}{chemical potential} $\mu$ is the change in
    internal energy associated with adding particles at fixed entropy and
    volume.
\end{itemize}
The \textcolor{blue}{second law}, in the form of the Clausius inequality,
states
\begin{equation}
    \boxed{\delta Q\leq T\,dS,}
\end{equation}
with equality for reversible processes. Combining the two laws for a
reversible process gives
\begin{equation}
    \boxed{dU=T\,dS-P\,dV+\mu\, dN,}
    \label{eq:combined-law}
\end{equation}
and therefore
\begin{equation}
    \boxed{
    \frac{\partial U}{\partial S}\bigg|_{V,N}=T, \qquad
    \frac{\partial U}{\partial V}\bigg|_{S,N}=-P, \qquad
    \frac{\partial U}{\partial N}\bigg|_{S,V}=\mu.}
    \label{eq:U-derivs}
\end{equation}
Rearranging Eq.~\eqref{eq:combined-law} gives the equivalent form
\begin{equation}
    \boxed{dS=\frac{1}{T}\,dU+\frac{P}{T}\,dV-\frac{\mu}{T}\,dN,}
\end{equation}
and therefore
\begin{equation}
    \boxed{
    \frac{\partial S}{\partial U}\bigg|_{V,N}=\frac{1}{T}, \qquad
    \frac{\partial S}{\partial V}\bigg|_{U,N}=\frac{P}{T}, \qquad
    \frac{\partial S}{\partial N}\bigg|_{U,V}=-\frac{\mu}{T}.}
    \label{eq:S-derivs}
\end{equation}

\subsection{Extensive and Intensive Variables}

\paragraph{Extensive and intensive variables and their conjugate pairs.}
\begin{itemize}
    \item A quantity $X$ is \textcolor{blue}{extensive} if, for two weakly
    coupled subsystems, $X_{\mathrm{total}}=X_1+X_2$. Equivalently, for
    $\lambda$ identical, non-interacting copies, $X_\lambda=\lambda X_1$.
    \item A quantity $Y$ is \textcolor{blue}{intensive} if it does not scale
    with the number of copies: $Y_\lambda=Y_1$.
    \item Energy, particle number, volume, and entropy are extensive, while
    temperature, pressure, and chemical potential are intensive.
\end{itemize}
This distinction explains why the Lagrange multipliers associated with
extensive constraints are intensive. For any observable $A$ with constrained
average $\langle A\rangle$, the multiplier is
\begin{equation*}
    \beta_A\equiv\frac{1}{k_B}\frac{\partial S}{\partial\langle A\rangle}.
\end{equation*}
Because $\langle A\rangle$ and $S$ are both extensive, their scaling factors
cancel in the derivative, so $\beta_A$ is intensive. Thermodynamic variables
therefore come in \textcolor{blue}{conjugate pairs}: an extensive variable
and its \textcolor{blue}{intensive conjugate}.

\paragraph{Additivity of entropy follows from statistical independence.}
\textcolor{red}{Here, the superscripts $(1)$ and $(2)$ also label the
distributions of the two systems, $p^{(1)}_i$ and $p^{(2)}_j$.}
The partition function factorizes because the joint distribution does:
$p_{ij}=p^{(1)}_i\,p^{(2)}_j$. The entropy is then additive:
\begin{align*}
    S_{1+2}
    = -\sum_{i,j} p^{(1)}_i p^{(2)}_j
    \ln\!\big(p^{(1)}_i p^{(2)}_j\big)
    = S_1+S_2.
\end{align*}
This requires statistical independence, not equilibrium, so the two systems
may have different multipliers $\beta_1$ and $\beta_2$:
\begin{equation}
    \boxed{S_{1+2}=k_B\left(\beta_1\langle U_1\rangle+\ln Z_1
    +\beta_2\langle U_2\rangle+\ln Z_2\right).}
\end{equation}
Only when $\beta_1=\beta_2=\beta$, that is, when the two systems are in
thermal equilibrium, does the joint system become canonical with a single
multiplier:
\begin{equation}
    \boxed{S_{1+2}=k_B\left[\beta\langle U_{1+2}\rangle+\ln Z_{1+2}\right],
    \qquad U_{1+2}=U_1+U_2.}
\end{equation}
Independence holds when the coupling between the systems is weak. For strong
coupling, the joint distribution does not factorize, and
$S_{1+2}=S_1+S_2-I$ with the mutual information $I\geq0$.

\subsection{Energy Fluctuation, Thermodynamic Limit, and Equivalence of Ensembles}

\paragraph{Energy fluctuation and the heat capacity.}
The variance of the internal energy, the \textcolor{blue}{energy
fluctuation} around the controlled mean $\langle U\rangle=\bar U$, can be
written in familiar thermodynamic quantities. By the chain rule,
\begin{align*}
    \mathrm{Var}(U)
    = -\frac{\partial\langle U\rangle}{\partial\beta}
    = -\frac{\partial\langle U\rangle}{\partial T}\frac{\partial T}{\partial\beta}.
\end{align*}
Since $T=1/(k_B\beta)$, we have $\partial T/\partial\beta=-k_BT^2$. By
definition, $\left.\partial\langle U\rangle/\partial T\right|_{V,N}\equiv C_V$
is the \textcolor{blue}{isochoric heat capacity}. It is computed at fixed
volume and particle number, which are the natural variables of the canonical
ensemble. Substituting,
\begin{equation}
    \boxed{\mathrm{Var}(U) = k_B T^2\, C_V.}
\end{equation}

\paragraph{Why fluctuations vanish in the thermodynamic limit.}
One might attribute the smallness of fluctuations to the tiny numerical
value of $k_B$ in SI units. This is not the reason: the value of $k_B$ only
reflects our choice of units. The genuine reason is that $C_V$ and
$\langle U\rangle$ are both extensive and scale linearly with the particle
number $N$. Hence
\begin{equation}
    \boxed{\frac{\sqrt{\mathrm{Var}(U)}}{\langle U\rangle}
    \sim\frac{\sqrt{N}}{N}=\frac{1}{\sqrt{N}},}
\end{equation}
which vanishes as $N\to\infty$. For a macroscopic system, relative energy
fluctuations are negligible because of the sheer size of $N$.
\textcolor{red}{For this reason, we write $U$ in place of $\langle U\rangle$
from here on, and likewise for other extensive variables.}
At large $N$, the canonical ensemble therefore behaves as the microcanonical
one: the fluctuation is negligible, and the relation between energy and
temperature is one-to-one.

\paragraph{Physical reading of the fluctuation formula.}
At higher temperature, $\beta$ is smaller, so the entropic cost per unit of
exchanged energy is smaller. The system then absorbs or releases energy more
freely, consistent with the larger fluctuation $k_BT^2C_V$.

\paragraph{Concavity of $S(U)$ and thermodynamic stability.}
Applying the concavity result to $A=U$ and using $\mathrm{Var}(U)=k_BT^2C_V$,
\begin{equation}
    \boxed{\frac{\partial^2 S}{\partial U^2}
    = -\frac{k_B}{\mathrm{Var}(U)} = -\frac{1}{C_V\,T^2}.}
\end{equation}
Concavity of $S(U)$ is therefore equivalent to $C_V>0$, the
\textcolor{blue}{thermodynamic stability} condition: adding energy to a
system must raise its temperature. This property is needed again when we
introduce Legendre transformations.

\subsection{Zeroth Law of Thermodynamics}

\paragraph{Two bodies exchanging energy.}
Consider two bodies, $(1)$ and $(2)$, in contact so that they can exchange
energy. Two points need to be stated precisely:
\begin{itemize}
    \item The combined system $(1)+(2)$ is genuinely isolated: nothing is
    exchanged with the outside. The total energy is a hard constraint,
    $U_1+U_2=U_{\text{tot}}$, so whatever one body releases, the other
    receives. This is a microcanonical condition on the pair as a whole.
    \item Each body has an entropy function $S_i(U_i,V_i,N_i)$, which is the
    canonical entropy derived earlier in this chapter. Each body carries the
    multiplier $\beta_i\equiv k_B^{-1}\,\partial S_i/\partial U_i$ attached to
    its own energy. The pair is microcanonical, while each body behaves
    canonically with the other body as its energy reservoir.
\end{itemize}

\paragraph{Equilibrium requires equal multipliers.}
The two bodies are weakly coupled, so the combined entropy is additive,
$S=S_1+S_2$. Equilibrium corresponds to the maximum of $S$ with respect to
how $U_{\text{tot}}$ is shared. Substituting $U_2=U_{\text{tot}}-U_1$ and
differentiating with respect to $U_1$,
\begin{align*}
    \frac{\partial S}{\partial U_1}
    = \frac{\partial S_1}{\partial U_1}
    +\frac{\partial S_2}{\partial U_2}\frac{\partial U_2}{\partial U_1}
    = \frac{\partial S_1}{\partial U_1}-\frac{\partial S_2}{\partial U_2}=0.
\end{align*}
Since $\beta_i=k_B^{-1}\,\partial S_i/\partial U_i$, this gives
\begin{equation}
    \boxed{\beta_1=\beta_2, \qquad\text{i.e.}\qquad T_1=T_2.}
\end{equation}

\paragraph{Direction of energy flow and the zeroth law.}
Suppose